\subsection{Qualitäts- und Wartbarkeitsanforderungen}
\label{subsec:qualitaets-wartbarkeitsanforderungen}

Die Anforderungen orientieren sich an überprüfbaren Merkmalen von
ISO/IEC~25010 \parencite{iso25010-2023}. Komponenten und Profile sollen klare
Verantwortlichkeiten besitzen, getrennt testbar und durch dokumentierte
Entscheidungen wartbar sein. Erweiterungen und Technologieaktualisierungen
sollen keine Kopie der gesamten Baseline erfordern, bleiben aber risikobehaftet.

Gleiche Eingaben und Konfigurationen sollen nachvollziehbare Generierungs- und
Build-Ergebnisse liefern. Dies stärkt die prüfbare Beziehung zwischen
Quellstand und Artefakt \parencite{lamb2022reproducible-builds}. Zielartefakte
sollen kontrolliert erzeugbar und Konfigurationsprioritäten verständlich sein;
serverseitige Geheimnisse bleiben von Client-Konfiguration getrennt.

Governance~v1 ergänzt funktional das Laden und Validieren einer versionierten
Policy, die Auswahl unterstützter Quelldateien, Code- und physische
Zeilenzählung, Scope-Metriken, Architekturregeln, Excludes, befristete
Exceptions sowie Text- und JSON-Reporting. Befunde sollen Pfad, optionale
Zeile und Symbol, Regel-ID und -name, Severity, Istwert, Schwelle, Erläuterung
und Handlungshinweis zuordnen. Nichtfunktional werden ein gemeinsamer lokaler
und CI-Einstieg, eine vom Dispatcher getrennte Quality-Logik, kontrollierte
Fehler bei ungültiger Konfiguration und die Verträglichkeit mit erzeugten
Profilen verlangt \parencite{kleiveist2026quality-policy}.

Tabelle~\ref{tab:quality-thresholds} gibt die implementierte reguläre
Klassifikation wieder. \emph{WARNING} und \emph{STRONG\_WARNING} informieren,
ohne allein den Exitcode zu verändern; \emph{ERROR} blockiert, sofern keine
greifende Exception vorliegt.

\begin{table}[H]
  \centering
  \scriptsize
  \renewcommand{\arraystretch}{1.08}
  \caption{Projektspezifische Schwellenwerte der Governance~v1}
  \label{tab:quality-thresholds}
  \begin{tabularx}{\textwidth}{@{}>{\raggedright\arraybackslash}X>{\centering\arraybackslash}p{0.13\textwidth}>{\centering\arraybackslash}p{0.15\textwidth}>{\centering\arraybackslash}p{0.18\textwidth}>{\centering\arraybackslash}p{0.13\textwidth}@{}}
    \toprule
    \textbf{Messgröße} & \textbf{PASS} & \textbf{WARNING} &
      \textbf{STRONG\_WARNING} & \textbf{ERROR} \\
    \midrule
    Code-LOC je Datei       & 0--600  & 601--750 & 751--900 & ab 901 \\
    physische Zeilen        & 0--1200 & ab 1201  & --       & -- \\
    Code-LOC je Funktion    & 0--50   & 51--80   & 81--120  & ab 121 \\
    Code-LOC je Klasse      & 0--300  & 301--500 & 501--700 & ab 701 \\
    zyklomatische Komplexität & 1--10 & 11--15   & 16--20   & ab 21 \\
    Verschachtelungstiefe   & 0--3    & 4        & 5        & ab 6 \\
    Parameter je Funktion   & 0--5    & 6--8     & 9--10    & ab 11 \\
    FastAPI-Handler, Code-LOC & 0--50 & --       & --       & ab 51 \\
    \bottomrule
  \end{tabularx}
  \smallskip
  {\footnotesize Quelle: Eigene Darstellung nach
  \textcite{kleiveist2026quality-config}.\par}
\end{table}

Die Grenze von 900 Codezeilen ist eine projektspezifische
Governance-Entscheidung, kein allgemeingültiger wissenschaftlicher
Clean-Code-Grenzwert: 900 LOC sind zulässig, liegen aber bereits in der starken
Warnzone; 901 LOC werden regulär als \texttt{CQ001 ERROR} klassifiziert. Ebenso
operationalisieren 120 Funktions-, 700 Klassenzeilen, Komplexität~20, fünf
Verschachtelungsebenen und zehn Parameter die Policy dieses Repository. Die
Quelle von McCabe begründet die zyklomatische Metrik, nicht den hier gewählten
Schwellenwert~20 \parencite{mccabe1976complexity}.

Das Dateilimit ist ein äußeres Sicherheitsnetz. Eine Datei mit 850 LOC kann
trotz zulässigem Hard-Limit problematisch sein; eine Datei mit 200 LOC kann
hohe Komplexität, ungünstige Abhängigkeiten oder vermischte Verantwortungen
besitzen. Die Metriken sind automatisierbare Indikatoren und ersetzen weder
fachliche Codeprüfung noch eine empirische Wartbarkeitsmessung. Die technische
Unterdrückbarkeit von \texttt{CQ001} wird deshalb gesondert als Abweichung in
Abschnitt~\ref{subsec:schwaechen-grenzen} bewertet.
\beginnerinput{workspace/statement/200_anforderungen/240}
